\chapter*{Convergence Tests for Positive-Term Series}
\addcontentsline{toc}{chapter}{Convergence Tests for Positive-Term Series}

\section*{The Integral Test}

Suppose $a_n=f(n)$ where $f$ is positive, continuous, and decreasing for $x\ge N$. Then the series and the corresponding improper integral behave the same way.

\begin{theorem}[Integral Test]
If $f(n)=a_n$ and $f$ is positive, continuous, and decreasing on $[N,\infty)$, then
$$\sum_{n=N}^{\infty} a_n$$
converges if and only if
$$\int_N^{\infty} f(x)\,dx$$
converges.
\end{theorem}

A classic application is the $p$-series
$$\sum_{n=1}^{\infty} \frac{1}{n^p}.$$ 
Using the integral test together with the improper integral $\int_1^{\infty}x^{-p}\,dx$, we obtain:

\begin{theorem}[$p$-series test]
The series
$$\sum_{n=1}^{\infty} \frac{1}{n^p}$$
converges if $p>1$ and diverges if $p\le 1$.
\end{theorem}

\section*{Direct Comparison Test}

When $0\le a_n\le b_n$, the larger series controls the smaller one.

\begin{theorem}[Direct Comparison]
Assume $0\le a_n\le b_n$ for all sufficiently large $n$.
\begin{itemize}
    \item If $\sum b_n$ converges, then $\sum a_n$ converges.
    \item If $\sum a_n$ diverges, then $\sum b_n$ diverges.
\end{itemize}
\end{theorem}

This test is most effective when we can compare with a geometric series or a $p$-series.

\section*{Limit Comparison Test}

Sometimes direct inequalities are awkward, but the terms have the same asymptotic size.

\begin{theorem}[Limit Comparison]
Let $a_n>0$ and $b_n>0$. If
$$\lim_{n\to\infty}\frac{a_n}{b_n} = c$$
for some constant $c$ with $0<c<\infty$, then either both series
$$\sum a_n, \qquad \sum b_n$$
converge, or both diverge.
\end{theorem}

\begin{remark}
Limit comparison is especially useful for rational expressions in $n$. In practice, one often compares using only the highest powers in the numerator and denominator.
\end{remark}